% Luc Destombe --- licence CC BY-NC-SA 4.0
% https://creativecommons.org/licenses/by-nc-sa/4.0/deed.fr
% Fichier autonome : compiler avec pdflatex, deux fois (signets, nombre de pages).
\documentclass[a4paper,10pt]{article}
\makeatletter
\newif\iflycee@niveaux
\lycee@niveauxtrue
\newif\iflycee@palatino
\lycee@palatinotrue

\RequirePackage[utf8]{inputenc}
\RequirePackage[T1]{fontenc}
\RequirePackage[french]{babel}
\RequirePackage{etoolbox}

\newcommand{\ShorthandsOff}{\@ifundefined{shorthandoff}{}{\shorthandoff{;:!?}}}
\newcommand{\ShorthandsOn}{\@ifundefined{shorthandon}{}{\shorthandon{;:!?}}}
\ShorthandsOff
\AtBeginDocument{\ShorthandsOn}
\AtBeginEnvironment{tikzpicture}{\ShorthandsOff}

\RequirePackage{geometry}
\RequirePackage{fancyhdr}
\RequirePackage{lastpage}
\RequirePackage{multicol}
\RequirePackage{array}
\RequirePackage{enumitem}
\RequirePackage{xcolor}
\RequirePackage{needspace}

\RequirePackage{amsmath,amssymb}
\RequirePackage{mathpazo}
\DeclareMathSizes{10.95}{10.95}{8}{5.5}
\begingroup\catcode`\_=\active \catcode`\^=\active
\gdef_#1{\sb{\scriptscriptstyle #1}}
\gdef^#1{\sp{\scriptscriptstyle\lycee@moinsepais #1}}
\endgroup
\newcommand\lycee@moins{\mathbin{%
  \setbox\z@\hbox{$\m@th\scriptscriptstyle\lycee@moinsorig$}%
  \dimen@=.55\wd\z@ \advance\dimen@-.5pt
  \hbox to.75\wd\z@{\hss$\m@th\scriptscriptstyle\vcenter{\hbox{%
    \tikz\draw[line width=.5pt,line cap=round](0,0)--(\the\dimen@,0);}}$\hss}}}
\begingroup\lccode`\~=`\- \lowercase{\endgroup\let~}\lycee@moins
\newcommand\lycee@moinsepais{\mathcode`\-="8000 }
\AtBeginDocument{\mathchardef\lycee@moinsorig=\mathcode`\-\relax}
\AtBeginDocument{\catcode`\_=12 \mathcode`\_="8000
  \catcode`\^=12 \mathcode`\^="8000 }
\newcommand\lycee@exposants{%
  \fontdimen13\textfont2=.48\fontdimen6\textfont2
  \fontdimen14\textfont2=.48\fontdimen6\textfont2
  \fontdimen15\textfont2=.48\fontdimen6\textfont2
  \fontdimen13\scriptfont2=.48\fontdimen6\scriptfont2
  \fontdimen14\scriptfont2=.48\fontdimen6\scriptfont2
  \fontdimen15\scriptfont2=.48\fontdimen6\scriptfont2}
\every@math@size\expandafter{\the\every@math@size\lycee@exposants}
\newlength{\retraitcalculs}
\setlength{\retraitcalculs}{2em}

\RequirePackage{tikz}
\usetikzlibrary{arrows.meta,calc,babel,shapes.misc}
\RequirePackage[most]{tcolorbox}
\tcbuselibrary{skins,breakable}

\definecolor{coulDef}{HTML}{1F4E79}
\definecolor{coulProp}{HTML}{7030A0}
\definecolor{coulRem}{HTML}{C55A11}
\definecolor{coulEx}{HTML}{2E7D32}
\definecolor{coulExo}{HTML}{B03A2E}
\definecolor{coulPart}{HTML}{1F4E79}
\definecolor{coulMeth}{HTML}{1B6E4F}
\definecolor{coulCourbe}{HTML}{0B5ED7}
\definecolor{coulCourbeB}{HTML}{C00000}
\definecolor{coulCourbeC}{HTML}{2E7D32}
\definecolor{coulGrille}{HTML}{8C8C8C}
\definecolor{coulRep}{HTML}{1B6E4F}
\definecolor{coulBar}{HTML}{2E7D32}

\newcommand{\R}{\mathbb{R}}
\newcommand{\Rs}{\mathbb{R}^{*}}
\renewcommand{\le}{\leqslant}
\renewcommand{\ge}{\geqslant}
\newcommand{\vect}[1]{\overrightarrow{#1}}
\newcommand{\norme}[1]{\left\lVert #1 \right\rVert}
\newcommand{\intervalle}[2]{\left[#1\,;#2\right]}
\RequirePackage{cancel}
\renewcommand{\CancelColor}{\color{coulExo}}
\newcommand{\fois}[1]{\times{\color{coulExo}#1}}
\newcommand{\barre}[1]{\cancel{#1}}

\tikzset{grille/.style={coulGrille,line width=0.4pt}}
\newcommand{\quadrillage}[4]{\draw[grille,step=1] (#1,#2) grid (#3,#4);}
\tikzset{
  croix/.style={cross out,draw,line width=0.6pt,inner sep=0pt,minimum size=4pt},
  croixdroite/.style={croix,rotate=45,line width=0.8pt},
}
\newcommand{\pt}[3]{\path (#1) node[croix]{} node[#2,font=\small,inner sep=2.5pt]{$#3$};}

\newcommand{\lycee@repere}[4]{%
  \draw[grille,step=1] (#1,#2) grid (#3,#4);
  \draw[-{Stealth[length=2mm]},thick] (#1,0)--({#3+0.4},0) node[below left=-1pt]{$x$};
  \draw[-{Stealth[length=2mm]},thick] (0,#2)--(0,{#4+0.4}) node[below left=-1pt]{$y$};
  \node[below left,font=\scriptsize,inner sep=1.5pt] at (0,0) {$0$};}
\newcommand{\lycee@gradx}[1]{\foreach \k/\l in {#1}{%
  \draw (\k,0.07)--(\k,-0.07) node[below,font=\scriptsize,inner sep=1.5pt]{$\l$};}}
\newcommand{\lycee@grady}[1]{\foreach \k/\l in {#1}{%
  \draw (0.07,\k)--(-0.07,\k) node[left,font=\scriptsize,inner sep=1.5pt]{$\l$};}}
\newcommand{\lycee@intff}[2]{\left[#1\,;#2\right]}
\newcommand{\lycee@intfo}[2]{\left[#1\,;#2\right[}
\newcommand{\lycee@intof}[2]{\left]#1\,;#2\right]}
\newcommand{\lycee@intoo}[2]{\left]#1\,;#2\right[}
\newcommand{\lycee@ens}[1]{\left\{#1\right\}}
\AtBeginDocument{%
  \providecommand{\repere}{\lycee@repere}%
  \providecommand{\gradx}{\lycee@gradx}%
  \providecommand{\grady}{\lycee@grady}%
  \providecommand{\Cf}{\mathcal{C}\sb{\scriptscriptstyle f}}%
  \providecommand{\Cg}{\mathcal{C}\sb{\scriptscriptstyle g}}%
  \providecommand{\intff}{\lycee@intff}%
  \providecommand{\intfo}{\lycee@intfo}%
  \providecommand{\intof}{\lycee@intof}%
  \providecommand{\intoo}{\lycee@intoo}%
  \providecommand{\ens}{\lycee@ens}}

\newlist{questions}{enumerate}{3}
\setlist[questions,1]{label=\arabic*\textdegree),leftmargin=*,itemsep=2pt,topsep=3pt}
\setlist[questions,2]{label=\alph*),leftmargin=*,itemsep=1pt,topsep=2pt}
\setlist[questions,3]{label=\roman*),leftmargin=*,itemsep=1pt,topsep=1pt}
\setlist[itemize]{label=\textbullet,leftmargin=*,itemsep=2pt,topsep=3pt}

\newcommand{\besoinplace}[1]{\par\penalty\@M\begingroup
  \dimen@=#1\relax\dimen@ii=\pagegoal\advance\dimen@ii-\pagetotal
  \ifdim\dimen@>\dimen@ii\ifdim\dimen@ii>\z@\vfil\fi\break\fi\endgroup}

\newcommand{\lycee@pied}{\thepage/\pageref{LastPage}}
\newcommand{\entete}[2]{%
  \pagestyle{fancy}\fancyhf{}%
  \fancyhead[L]{\footnotesize\color{coulPart}\textbf{#1}}%
  \fancyhead[R]{\footnotesize\color{coulPart}\textbf{#2}}%
  \fancyfoot[C]{\footnotesize\lycee@pied}%
  \renewcommand{\headrulewidth}{0.4pt}%
  \renewcommand{\headrule}{\hbox to\headwidth{%
    \color{coulPart!40}\leaders\hrule height \headrulewidth\hfill}}}

\RequirePackage{ccicons}
\newcommand{\lycee@licence}{{\scriptsize\color{gray}%
  \copyright\ Luc Destombe --- \ccbyncsa\ CC BY-NC-SA 4.0}}
\newif\iflycee@licence \lycee@licencetrue
\newcommand{\sanslicence}{\lycee@licencefalse}
\AtBeginDocument{\iflycee@licence\AddToHookNext{shipout/foreground}{%
  \put(0,0){\rlap{\kern\dimexpr1in+\hoffset+\oddsidemargin+\textwidth\relax
    \llap{\raisebox{-\dimexpr1in+\voffset+\topmargin+\headheight+\headsep
      +\textheight+\footskip\relax}{\lycee@licence}}}}}\fi}

\newcommand{\formulesserrees}{%
  \setlength{\abovedisplayskip}{3pt}\setlength{\belowdisplayskip}{3pt}%
  \setlength{\abovedisplayshortskip}{2pt}\setlength{\belowdisplayshortskip}{3pt}}

\setlength{\parindent}{0pt}

\RequirePackage[implicit=false,hidelinks,pdfencoding=auto,psdextra,bookmarksopen,
  bookmarksopenlevel=1,pdfcreator={},pdfproducer={}]{hyperref}
\RequirePackage{bookmark}
\hypersetup{pdfauthor={Luc Destombe},pdfkeywords={CC BY-NC-SA 4.0}}
\newcounter{lycee@signet}
\newcommand{\signet}[3][]{\stepcounter{lycee@signet}%
  \edef\lycee@dest{\if\relax\detokenize{#1}\relax
    signet.\arabic{lycee@signet}\else#1\fi}%
  \hypertarget{\lycee@dest}{}\bookmark[level=#2,dest=\lycee@dest]{#3}}
\pdfstringdefDisableCommands{%
  \def\R{R}\def\vect#1{#1}\def\Cf{Cf}\def\Cg{Cg}\def\fois#1{×#1}%
  \def\barre#1{#1}\def\intervalle#1#2{[#1;#2]}\def\intff#1#2{[#1;#2]}%
  \def\textsuperscript#1{#1}\def\dfrac#1#2{#1/#2}\def\frac#1#2{#1/#2}%
  \def\vec#1{#1⃗}}%
\RequirePackage[official]{eurosym}

\geometry{top=0.8cm,bottom=1.3cm,left=1cm,right=1cm,footskip=0.6cm}

\pagestyle{fancy}
\fancyhf{}
\renewcommand{\headrulewidth}{0pt}
\fancyfoot[C]{\footnotesize Page \thepage{} / \pageref{LastPage}}

\newcommand{\titrefiche}[2]{%
  \begin{center}
    \Large\textbf{#1}\\[2pt]
    \large\textbf{#2}\\[2pt]
  \end{center}
  \vspace{0.10cm}}

\setlength{\columnseprule}{0.4pt}
\setlength{\columnsep}{15pt}
\renewcommand{\columnseprulecolor}{\color{black!45}}
\setlength{\multicolsep}{2pt}

\newcounter{partiefiche}
\renewcommand{\thepartiefiche}{\Roman{partiefiche}}
\newif\ifapresbandeau
\newcommand{\lycee@bandeau}[1]{%
  \par\penalty-600
  \vspace{0.08cm}%
  \noindent\colorbox{coulPart}{%
    \parbox{\dimexpr\linewidth-2\fboxsep\relax}{%
      \color{white}\bfseries\raggedright\strut#1\strut}}%
  \nopagebreak\par\nopagebreak
  \vspace{0.06cm}\nopagebreak
  \apresbandeautrue}
\newcommand{\partiefiche}[1]{%
  \refstepcounter{partiefiche}%
  \lycee@bandeau{\signet{1}{\thepartiefiche{} --- #1}\thepartiefiche{} --- #1}}
\newcommand{\partiefichesansnum}[1]{\lycee@bandeau{\signet{1}{#1}#1}}

\newcounter{exo}
\newlength{\espaceexo}\setlength{\espaceexo}{7pt}
\newlength{\espacetitre}\setlength{\espacetitre}{2pt}
\newcommand{\exo}[1][\relax]{%
  \par
  \ifapresbandeau\apresbandeaufalse\else\penalty-150\addvspace{\espaceexo}\fi
  \refstepcounter{exo}%
  \noindent\signet[exercice-\arabic{exo}]{2}{Exercice \arabic{exo}}%
  \textbf{\color{coulExo}\underline{Exercice \theexo}}%
  \ifx\relax#1\relax\else\textbf{ : #1}\fi
  \par\nopagebreak\vspace{\espacetitre}\nopagebreak
  \ignorespaces}

\setlist[enumerate]{nosep,leftmargin=*,itemsep=2pt,topsep=2pt}
\setlist[itemize]{nosep,leftmargin=*,topsep=2pt}
\newcommand{\choix}[3]{\par\hspace*{1em}%
  \makebox[0.3\linewidth][l]{a) #1}\makebox[0.3\linewidth][l]{b) #2}c) #3}
\newcommand{\deux}[2]{\par\hspace*{0.6em}%
  \makebox[0.48\linewidth][l]{#1}#2}
\newcommand{\trois}[3]{\par\hspace*{0.6em}%
  \makebox[0.32\linewidth][l]{#1}\makebox[0.32\linewidth][l]{#2}#3}

  \renewcommand{\exo}[1][\relax]{%
    \ifapresbandeau\apresbandeaufalse\else\par\penalty-150\fi
    \refstepcounter{exo}%
    \vspace{0.05cm}%
    \noindent\signet[exercice-\arabic{exo}]{2}{Exercice \arabic{exo}}%
    \textbf{\color{coulExo}\underline{Exercice \theexo}}%
    \ifx\relax#1\relax\else\textbf{ : #1}\fi%
    \par\nopagebreak
    \ignorespaces}
  \newcommand{\rappel}[1]{%
    \par\vspace{0.06cm}\nopagebreak
    \noindent\fcolorbox{black!35}{black!5}{%
      \parbox{\dimexpr\linewidth-2\fboxsep-2\fboxrule\relax}{%
        \small\textbf{Rappel.} #1}}%
    \par\vspace{0.06cm}\nopagebreak}
  \setlist[enumerate]{nosep,leftmargin=*}
  \setlist[itemize]{nosep,leftmargin=*}
  \setlength{\multicolsep}{3pt plus 1pt minus 1pt}
  \setlength{\premulticols}{10pt}
  \setlength{\postmulticols}{10pt}
  \newlist{expr}{enumerate}{1}
  \setlist[expr]{label={\Alph*\ $=$},nosep,leftmargin=*,itemsep=0pt}

\renewenvironment{center}{\par\addvspace{3pt}\centering}{\par\addvspace{3pt}}
\formulesserrees
\AtBeginDocument{\formulesserrees}
\clubpenalty=10000
\widowpenalty=10000
\displaywidowpenalty=10000
\brokenpenalty=10000
\setlength{\parskip}{0pt}

\RequirePackage{ragged2e}
\setlength{\RaggedRightRightskip}{0pt plus 1fil}
\AtBeginDocument{\RaggedRight\binoppenalty=10000 \relpenalty=10000 }
\g@addto@macro\@parboxrestore{\RaggedRight}
\makeatother
\geometry{top=0.8cm,bottom=1.3cm,left=1cm,right=1cm,footskip=0.7cm,headheight=12pt}

\raggedcolumns
\newcommand{\qcm}[4]{\par\makebox[0.24\linewidth][l]{a) #1}%
  \makebox[0.24\linewidth][l]{b) #2}\makebox[0.24\linewidth][l]{c) #3}d) #4}

\begin{document}

\titrefiche{1\textsuperscript{re} mathématiques spécifiques --- Fonctions affines et pourcentages}%
  {Fiche d'exercices du chapitre}

\begin{multicols}{2}

\partiefiche{Appliquer un pourcentage}

\exo[Calcul mental]
Calculer de tête.
\trois{a) $10\,\%$ de $420$}{b) $25\,\%$ de $60$}{c) $50\,\%$ de $48$}
\trois{d) $1\,\%$ de $900$}{e) $5\,\%$ de $80$}{f) $75\,\%$ de $20$}
\trois{g) $20\,\%$ de $35$}{h) $30\,\%$ de $250$}{i) $40\,\%$ de $15$}

\exo[Trois écritures d'un même nombre]
Recopier et compléter le tableau.
\begin{center}
\renewcommand{\arraystretch}{1.9}
\begin{tabular}{|l|c|c|c|c|c|}
\hline
Fraction & $\dfrac{3}{5}$ & & & $\dfrac{1}{4}$ & \\
\hline
Décimal & & $0{,}09$ & & & $1{,}5$\\
\hline
Pourcentage & & & $45\,\%$ & & \\
\hline
\end{tabular}
\end{center}

\exo[Retrouver le total]
Pour chaque situation, calculer l'effectif total.
\begin{enumerate}[label=\alph*)]
  \item $30$ chevaux représentent $5\,\%$ des animaux d'une ferme.
  \item $9$ élèves représentent $30\,\%$ d'une classe.
  \item $8$~€ représentent $10\,\%$ du prix d'un jeu.
  \item $60$ adhérents représentent $75\,\%$ d'un club.
\end{enumerate}

\exo[Augmenter ou diminuer]
Calculer chaque nouvelle valeur, en calculant d'abord le montant de la hausse ou de la
baisse.
\begin{enumerate}[label=\alph*)]
  \item Un article à $60$~€ augmente de $15\,\%$.
  \item Un ordinateur à $900$~€ est soldé à $-20\,\%$.
  \item Un salaire de $2\,000$~€ augmente de $3\,\%$.
  \item Une commune de $6\,000$ habitants perd $5\,\%$ de sa population.
\end{enumerate}

\exo[Prix hors taxe, prix TTC]
Une trottinette coûte $350$~€ hors taxe (HT). On lui ajoute la TVA, de $20\,\%$.
\begin{enumerate}[label=\alph*)]
  \item Calculer le montant de la TVA, puis le prix toutes taxes comprises (TTC).
  \item En soldes, le prix TTC baisse de $10\,\%$. Calculer le prix soldé.
\end{enumerate}

\partiefiche{Calculer un pourcentage}

\exo[Proportions]
Écrire chaque proportion en pourcentage.
\deux{a) $12$ élèves sur $40$}{b) $9$ admis sur $20$ candidats}
\deux{c) $33$ garçons sur $150$ élèves}{d) $80$ clients sur $500$}

\exo[Deux enquêtes]
Répondre en justifiant par un calcul.
\begin{enumerate}[label=\alph*)]
  \item Sur $250$ élèves, $140$ sont favorables à un projet de jardin. Le projet
        exige qu'au moins $60\,\%$ d'entre eux le soient. Est-ce le cas ?
  \item Sur $400$ clients d'une agence de voyages, $340$ sont satisfaits. Une
        publicité affirme que $85\,\%$ d'entre eux le sont. Est-ce exact ?
\end{enumerate}

\exo[Taux d'évolution]
Calculer le taux d'évolution de chaque grandeur, en pourcentage.
\begin{enumerate}[label=\alph*)]
  \item Un prix passe de $150$~€ à $180$~€.
  \item Un stock passe de $400$~kg à $340$~kg.
  \item Un loyer passe de $800$~€ à $720$~€.
  \item Un club passe de $60$ à $75$ adhérents.
\end{enumerate}

\exo[Plus ou moins de\ldots{} ?]
Répondre sans calculer le taux d'évolution exact.
\begin{enumerate}[label=\alph*)]
  \item Un club passe de $700$ à $620$ adhérents. A-t-il perdu plus de $10\,\%$ de ses
        adhérents ?
  \item Une forêt passe de $1\,500$ à $1\,620$ arbres. A-t-elle gagné plus de
        $10\,\%$ d'arbres ?
  \item Une ville passe de $3\,000$ à $2\,880$ habitants. A-t-elle perdu plus de
        $5\,\%$ de ses habitants ?
\end{enumerate}

\exo[Au tableur]
Une feuille de calcul donne le nombre d'entrées d'une piscine.
\begin{center}
\renewcommand{\arraystretch}{1.25}
\small
\begin{tabular}{|c|l|c|c|c|c|}
\hline
 & \multicolumn{1}{c|}{A} & B & C & D & E\\
\hline
1 & Année & 2022 & 2023 & 2024 & 2025\\
\hline
2 & Entrées & $5\,000$ & $5\,500$ & $4\,950$ & $5\,445$\\
\hline
3 & Taux d'évolution & \large$\times$ & & & \\
\hline
\end{tabular}
\end{center}
\rappel{Dans une formule, \texttt{B2} désigne le nombre de la cellule B2. En recopiant
une formule vers la droite, \texttt{B2} devient \texttt{C2}, puis \texttt{D2}\ldots}
\begin{enumerate}[label=\alph*)]
  \item Calculer le taux d'évolution du nombre d'entrées entre 2022 et 2023.
  \item Parmi ces formules, laquelle saisir en C3 pour obtenir, par recopie vers la
        droite, les taux d'évolution d'une année sur l'autre ?
        \par\hspace*{0.6em}\makebox[0.45\linewidth][l]{\texttt{=(C2-B2)/C2}}\texttt{=C2-B2/B2}
        \par\hspace*{0.6em}\makebox[0.45\linewidth][l]{\texttt{=(C2-B2)/B2}}\texttt{=(B2-C2)/B2}
  \item Quelle formule contient alors la cellule D3 ?
\end{enumerate}

\partiefiche{Coefficient multiplicateur}

\exo[Du taux au coefficient]
Donner le coefficient multiplicateur de chaque évolution.
\trois{a) $+6\,\%$}{b) $-6\,\%$}{c) $+30\,\%$}
\trois{d) $-30\,\%$}{e) $+65\,\%$}{f) $-12\,\%$}
\trois{g) $+4\,\%$}{h) $-45\,\%$}{i) $+200\,\%$}

\exo[Du coefficient au taux]
Donner l'évolution, en pourcentage, associée à chaque coefficient.
\trois{a) $\times 1{,}09$}{b) $\times 0{,}91$}{c) $\times\frac{4}{5}$}
\trois{d) $\times 0{,}7$}{e) $\times 3$}{f) $\times 0{,}98$}

\exo[Appliquer une évolution]
Calculer chaque nouvelle valeur en une seule multiplication.
\begin{enumerate}[label=\alph*)]
  \item $300$~€ augmentés de $15\,\%$.
  \item $600$~€ diminués de $35\,\%$.
  \item $50$~€ augmentés de $120\,\%$.
  \item $2\,000$ habitants diminués de $1{,}5\,\%$.
\end{enumerate}

\exo[Le bon calcul]
Pour chaque question, une seule réponse est exacte.
\begin{enumerate}[label=\arabic*.]
  \item Un article à $500$~€ baisse de $12\,\%$. Son nouveau prix est :
        \qcm{$500-0{,}12$}{$500\times 0{,}88$}{$500\times 1{,}12$}{$500\times 0{,}12$}
  \item Un article à $150$~€ augmente de $30\,\%$. Son nouveau prix est :
        \qcm{$150+0{,}3$}{$150\times\dfrac{30}{100}$}{$150\times 1{,}3$}{$150\times 0{,}7$}
\end{enumerate}

\exo[Retrouver un taux avec le coefficient]
Calculer le coefficient multiplicateur, puis le taux d'évolution.
\begin{enumerate}[label=\alph*)]
  \item Un prix passe de $80$~€ à $100$~€.
  \item Un stock passe de $500$~kg à $400$~kg.
  \item Une production passe de $20$ à $60$ tonnes.
\end{enumerate}

\partiefiche{Évolutions successives, évolution réciproque}

\exo[Taux global]
Calculer le taux d'évolution global.
\begin{enumerate}[label=\alph*)]
  \item Une hausse de $50\,\%$, puis une hausse de $20\,\%$.
  \item Une baisse de $30\,\%$, puis une hausse de $30\,\%$.
  \item Une hausse de $100\,\%$, puis une baisse de $50\,\%$.
  \item Deux baisses successives de $20\,\%$.
  \item Une hausse de $60\,\%$, puis une baisse de $25\,\%$.
\end{enumerate}

\exo[Le prix de l'électricité]
Un journaliste annonce : « Une hausse de $10\,\%$ en février, puis une autre de
$10\,\%$ en août, augmenteront le prix de l'électricité de $20\,\%$ sur l'année. »
A-t-il raison ?

\exo[Vrai ou faux ?]
Justifier chaque réponse.
\begin{enumerate}[label=\alph*)]
  \item Une hausse de $50\,\%$ suivie d'une baisse de $50\,\%$ ne change rien.
  \item Deux remises successives de $20\,\%$ et de $30\,\%$ donnent le même prix,
        quel que soit leur ordre.
  \item Deux baisses successives de $10\,\%$ font une baisse de $19\,\%$.
\end{enumerate}

\exo[La valeur d'un scooter]
Un scooter électrique neuf coûte $4\,000$~€. Il perd $25\,\%$ de sa valeur la
première année, puis $20\,\%$ la deuxième année.
\begin{enumerate}[label=\alph*)]
  \item Calculer sa valeur au bout d'un an, puis de deux ans.
  \item Calculer le taux d'évolution global sur les deux ans.
\end{enumerate}

\exo[Une même évolution répétée]
Donner chaque fois le taux d'évolution global.
\begin{enumerate}[label=\alph*)]
  \item Une entreprise augmente sa production de $3\,\%$ par an. Par quel nombre la
        production est-elle multipliée en $3$ ans ? On donne
        $1{,}03^{3}\approx 1{,}09$.
  \item Les prix augmentent de $1\,\%$ par mois pendant un an. On donne
        $1{,}01^{12}\approx 1{,}13$. Quel est le taux d'évolution annuel ?
\end{enumerate}

\exo[Se compensent-elles ?]
Pour chaque cas, dire si la seconde évolution ramène à la valeur de départ, en
calculant le produit des coefficients multiplicateurs.
\deux{a) $-20\,\%$, puis $+25\,\%$}{b) $+50\,\%$, puis $-50\,\%$}
\deux{c) $-50\,\%$, puis $+100\,\%$}{d) $+300\,\%$, puis $-75\,\%$}
\deux{e) $-20\,\%$, puis $+20\,\%$}{f) $+10\,\%$, puis $-10\,\%$}

\exo[Quelle évolution compense ?]
Pour chaque question, une seule réponse est exacte : tester chaque réponse.
\begin{enumerate}[label=\arabic*.]
  \item Un prix est multiplié par $4$. Pour revenir au prix de départ, il faut une
        baisse de :
        \qcm{$400\,\%$}{$75\,\%$}{$25\,\%$}{$300\,\%$}
  \item Un prix augmente de $25\,\%$. Pour revenir au prix de départ, il faut une
        baisse de :
        \qcm{$25\,\%$}{$75\,\%$}{$20\,\%$}{$80\,\%$}
  \item Une population baisse de $50\,\%$. Pour revenir à son niveau, il faut une
        hausse de :
        \qcm{$100\,\%$}{$50\,\%$}{$200\,\%$}{$150\,\%$}
\end{enumerate}

\exo[Un verger]
Un verger compte $2\,000$ arbres. Une maladie en fait disparaître $20\,\%$ en un an.
\begin{enumerate}[label=\alph*)]
  \item Combien d'arbres reste-t-il ?
  \item Le propriétaire replante l'année suivante $20\,\%$ du nombre d'arbres restants.
        Retrouve-t-il $2\,000$ arbres ?
  \item Montrer qu'il suffirait de replanter $25\,\%$ du nombre d'arbres restants.
\end{enumerate}

\partiefiche{Fonctions affines : image, antécédent}

\exo[Affine ou non ?]
Dire si chaque fonction est affine ; si oui, donner $a$ et $b$.
\trois{a) $f(x)=5x+1$}{b) $g(x)=x^{2}-3$}{c) $h(x)=2-3x$}
\trois{d) $k(x)=\dfrac{2}{x}$}{e) $m(x)=1{,}5x$}{f) $n(x)=-4$}

\exo[Images]
Soit $f(x)=3x-4$ et $g(x)=-4x+6$.
\begin{enumerate}[label=\alph*)]
  \item Calculer $f(0)$, $f(2)$, $f(-1)$ et $f(0{,}5)$.
  \item Calculer $g(0)$, $g(2)$, $g(-3)$ et $g(1{,}5)$.
\end{enumerate}

\exo[Antécédents]
Calculer chaque antécédent par une équation.
\begin{enumerate}[label=\alph*)]
  \item $f(x)=3x+9$ : calculer l'antécédent de $15$, de $0$, puis de $-6$.
  \item $g(x)=-2x+8$ : calculer l'antécédent de $0$, puis de $14$.
  \item $h(x)=0{,}5x-2$ : résoudre $h(x)=3$.
\end{enumerate}

\exo[Deux tarifs]
Répondre à chaque question par un calcul.
\begin{enumerate}[label=\alph*)]
  \item Un dépanneur facture $40$~€ de déplacement et $1{,}50$~€ par km. Exprimer le
        prix $T(x)$ pour $x$~km. Quelle distance a-t-il parcourue pour $70$~€ ?
  \item Une salle de réception se loue $200$~€, plus $15$~€ par invité. Exprimer le
        prix $S(x)$ pour $x$ invités. Calculer $S(40)$. Combien d'invités pour
        $650$~€ ?
\end{enumerate}

\partiefiche{Trouver une fonction affine}

\exo[Mettre en équation]
Exprimer chaque grandeur par une fonction affine de $x$.
\begin{enumerate}[label=\alph*)]
  \item Une location de kayak coûte $8$~€, plus $6$~€ par heure. $L(x)$ est le prix
        pour $x$ heures.
  \item Un réservoir de $60$~L consomme $5$~L tous les $100$~km. $R(x)$ est la quantité
        restante après $x$~km.
  \item Une vendeuse touche $1\,400$~€ par mois, plus $4\,\%$ de ses ventes. $S(x)$ est
        son salaire pour $x$~€ de ventes.
\end{enumerate}

\exo[Coefficient directeur]
Calculer le coefficient directeur de la droite $(AB)$.
\deux{a) $A(1\,;2)$ et $B(4\,;11)$}{b) $A(0\,;6)$ et $B(2\,;2)$}
\deux{c) $A(-2\,;3)$ et $B(-4\,;7)$}{d) $A(-3\,;-2)$ et $B(1\,;6)$}

\exo[Déterminer une fonction affine]
Pour chaque cas, déterminer la fonction affine $f$ dont la droite passe par $A$ et $B$.
\deux{a) $A(1\,;5)$ et $B(3\,;11)$}{b) $A(0\,;-4)$ et $B(2\,;2)$}
\deux{c) $A(-2\,;3)$ et $B(2\,;3)$}{d) $A(-4\,;-5)$ et $B(2\,;4)$}

\exo[Le dioxyde de carbone dans l'air]
La concentration de dioxyde de carbone dans l'air était d'environ $370$~ppm en 2000 et
$390$~ppm en 2010. On la modélise par une fonction affine $f$ du nombre $x$ d'années
écoulées depuis 2000.
\begin{enumerate}[label=\alph*)]
  \item Donner $f(0)$ et $f(10)$, puis déterminer $f(x)$.
  \item Interpréter le coefficient directeur.
  \item Selon ce modèle, quelle sera la concentration en 2030 ?
  \item En quelle année atteindra-t-elle $450$~ppm ?
\end{enumerate}

\partiefiche{Tracer une fonction affine}

\exo[Sens de variation]
Donner le sens de variation de chaque fonction affine.
\trois{a) $f(x)=4-x$}{b) $g(x)=3x+1$}{c) $k(x)=-2$}
\trois{d) $p(x)=-5x$}{e) $h(x)=-0{,}2x+3$}{}
\par\hspace*{0.6em}f) $m(x)=0{,}3x-1$

\exo[Tracer des droites]
Dans un même repère orthonormé, tracer les droites qui représentent :
\trois{$f(x)=x+1$}{$g(x)=-3x+2$}{$h(x)=0{,}5x-2$}

\exo[Location de paddles]
Un paddle se loue $4$~€, plus $4$~€ par heure. On note $L(x)$ le prix pour $x$
heures, avec $x$ entre $0$ et $8$.
\begin{enumerate}[label=\alph*)]
  \item Exprimer $L(x)$ en fonction de $x$.
  \item Calculer $L(0)$ et $L(8)$.
  \item Tracer la droite qui représente $L$ (en abscisse, $1$~cm pour $1$ heure ; en
        ordonnée, $1$~cm pour $5$~€).
  \item Lire graphiquement la durée de location pour $24$~€, puis vérifier par le
        calcul.
\end{enumerate}

\exo[Proportionnalité et fonctions linéaires]
Répondre à chaque question par un calcul.
\begin{enumerate}[label=\alph*)]
  \item Une boisson contient $10$~g de sucre pour $100$~mL. Exprimer la masse de sucre
        $f(x)$ dans $x$~mL de boisson, puis calculer $f(330)$.
  \item Un pouce vaut environ $2{,}5$~cm. Convertir $40$ pouces en centimètres, puis
        $75$~cm en pouces.
  \item Parmi $f(x)=-2x$ ; $g(x)=2x+3$ ; $h(x)=\dfrac{x}{4}$ et $k(x)=-1$, lesquelles
        sont linéaires ?
  \item La droite d'une fonction linéaire $\ell$ passe par le point $(5\,;15)$.
        Déterminer $\ell(x)$.
\end{enumerate}

\partiefiche{Lire l'équation d'une droite}

\exo[Lire une équation réduite]
\begin{minipage}[t]{0.36\linewidth}\raggedright\vspace{0pt}
Donner l'équation réduite de chaque droite, puis le sens de variation de la fonction
qu'elle représente.
\end{minipage}\hfill\begin{minipage}[t]{0.62\linewidth}\centering\vspace{0pt}
\begin{tikzpicture}[x=0.4cm,y=0.4cm]
  \quadrillage{-4}{-4}{5}{5}
  \draw[->,>=Stealth,black!70] (-4,0) -- (5.5,0) node[below,font=\tiny] {$x$};
  \draw[->,>=Stealth,black!70] (0,-4) -- (0,5.5) node[left,font=\tiny] {$y$};
  \foreach \i in {-4,...,-1,1,2,...,4} {\node[below,font=\tiny] at (\i,0) {$\i$};}
  \foreach \i in {-4,...,-1,1,2,...,4} {\node[left,font=\tiny] at (0,\i) {$\i$};}
  \node[below left,font=\tiny] at (0,0) {$0$};
  \begin{scope}
    \clip (-4,-4) rectangle (5,5);
    \draw[coulCourbe,thick] (-1,-4) -- (2,5);
    \draw[coulCourbeB,thick] (-3,5) -- (5,-3);
    \draw[coulCourbeC,thick] (-4,-0.333) -- (5,2.667);
    \draw[black!70,thick] (-4,-2) -- (5,-2);
  \end{scope}
  \node[coulCourbe,right,font=\small] at (2,4.4) {$(d_{1})$};
  \node[coulCourbeB,left,font=\small] at (-2.6,4.4) {$(d_{2})$};
  \node[coulCourbeC,above,font=\small] at (4.2,2.4) {$(d_{3})$};
  \node[black!70,above,font=\small] at (4.2,-2) {$(d_{4})$};
\end{tikzpicture}
\end{minipage}

\exo[Points d'une droite]
Soit $(d)$ la droite d'équation réduite $y=-2x+5$.
\begin{enumerate}[label=\alph*)]
  \item Les points $A(1\,;3)$, $B(3\,;1)$ et $C(-2\,;9)$ sont-ils sur $(d)$ ?
  \item Calculer l'ordonnée du point de $(d)$ d'abscisse $4$.
  \item Calculer l'abscisse du point de $(d)$ d'ordonnée $-7$.
\end{enumerate}

\exo[Équation réduite sans graphique]
Pour chaque question, une seule réponse est exacte.
\begin{enumerate}[label=\arabic*.]
  \item Une droite coupe l'axe des ordonnées en $2$ ; pour un déplacement horizontal
        de $3$, son déplacement vertical est $+1$. Son équation réduite est :
        \deux{a) $y=3x+2$}{b) $y=\dfrac{1}{3}x+2$}
        \deux{c) $y=x+\dfrac{1}{3}$}{d) $y=\dfrac{1}{3}x-2$}
  \item La droite qui passe par $(0\,;-3)$ et $(2\,;1)$ a pour équation réduite :
        \deux{a) $y=-3x+2$}{b) $y=0{,}5x-3$}
        \deux{c) $y=2x-3$}{d) $y=-2x-3$}
  \item La droite d'équation $y=4x-1$ :
        \par a) monte \quad b) descend \quad c) est horizontale
\end{enumerate}

\partiefiche{Signe d'une fonction affine}

\exo[Tableaux de signes]
Pour chaque fonction, calculer la valeur qui l'annule, puis dresser son tableau de
signes.
\deux{a) $f(x)=3x-9$}{b) $g(x)=-2x+8$}
\deux{c) $h(x)=4x+12$}{d) $k(x)=-0{,}5x+3$}

\exo[Positif ou négatif ?]
Pour chaque question, une seule réponse est exacte.
\begin{enumerate}[label=\arabic*.]
  \item $f(x)=-5x+2$ est positif ou nul pour :
        \par\hspace*{0.6em}\makebox[0.45\linewidth][l]{a) $x\le 0{,}4$}b) $x\ge 0{,}4$
        \par\hspace*{0.6em}\makebox[0.45\linewidth][l]{c) $x\le -0{,}4$}d) $x\ge -0{,}4$
  \item $3x+12$ est strictement négatif pour :
        \qcm{$x>-4$}{$x<4$}{$x<-4$}{$x>4$}
\end{enumerate}

\exo[Un bénéfice]
Un food-truck vend $x$ repas par jour, avec $x$ entre $0$ et $80$. Son bénéfice, en
euros, est $B(x)=5x-150$.
\begin{enumerate}[label=\alph*)]
  \item Calculer $B(0)$ et interpréter ce résultat.
  \item Résoudre $B(x)=0$.
  \item Dresser le tableau de signes de $B(x)$ pour $x$ entre $0$ et $80$.
  \item À partir de combien de repas vendus le bénéfice est-il positif ?
\end{enumerate}

\partiefiche{Comparer deux fonctions affines}

\exo[Équations]
Résoudre chaque équation.
\deux{a) $6x+5=4x+11$}{b) $3x-2=x+8$}
\deux{c) $2x+3=5x-9$}{d) $4x+1=-2x+4$}

\exo[Deux offres d'emploi]
Une commerciale hésite entre deux sociétés. Société A : $1\,200$~€ par mois, plus
$10\,\%$ des ventes. Société B : $800$~€ par mois, plus $20\,\%$ des ventes. On note
$x$ le montant des ventes du mois, en euros.
\begin{enumerate}[label=\alph*)]
  \item Exprimer les salaires $f(x)$ et $g(x)$ proposés par A et par B.
  \item Pour quel montant de ventes les deux salaires sont-ils égaux ?
  \item Elle pense vendre pour $5\,000$~€ par mois. Quelle société choisir ?
\end{enumerate}

\exo[Deux formules de piscine]
Formule A : $5$~€ l'entrée. Formule B : abonnement de $24$~€, puis $2$~€ l'entrée. Le
graphique donne le prix payé selon le nombre $x$ d'entrées.
\begin{center}
\begin{tikzpicture}[x=0.45cm,y=0.05cm]
  \draw[grille,xstep=1,ystep=10] (0,0) grid (12,60);
  \draw[->,>=Stealth,black!70] (0,0) -- (12.7,0) node[below,font=\tiny] {$x$};
  \draw[->,>=Stealth,black!70] (0,0) -- (0,66) node[left,font=\tiny] {€};
  \foreach \i in {1,...,12} {\node[below,font=\tiny] at (\i,0) {$\i$};}
  \foreach \v in {10,20,...,60} {\node[left,font=\tiny] at (0,\v) {$\v$};}
  \node[below left,font=\tiny] at (0,0) {$0$};
  \draw[coulCourbe,thick] (0,0) -- (12,60);
  \draw[coulCourbeB,thick] (0,24) -- (12,48);
  \node[coulCourbe,left,font=\small] at (10.6,54) {A};
  \node[coulCourbeB,below,font=\small] at (11.4,45) {B};
\end{tikzpicture}
\end{center}
\begin{enumerate}[label=\alph*)]
  \item Exprimer les prix $f(x)$ et $g(x)$ des formules A et B.
  \item Lire les coordonnées du point d'intersection des deux droites.
  \item Retrouver ce résultat par le calcul.
  \item Quelle formule choisir selon le nombre d'entrées ?
\end{enumerate}

\partiefichesansnum{Exercices type épreuve anticipée}

\exo[Automatismes (QCM)]
Pour chaque question, une seule réponse est exacte.
\begin{enumerate}[label=\arabic*.]
  \item $30\,\%$ de $70$ est égal à :
        \qcm{$21$}{$210$}{$40$}{$2{,}1$}
  \item Un article à $100$~€ augmente de $20\,\%$, puis baisse de $20\,\%$. Il coûte :
        \qcm{$100$~€}{$96$~€}{$104$~€}{$80$~€}
  \item Lequel de ces nombres est inférieur à $0{,}2$ ?
        \qcm{$\dfrac{1}{4}$}{$25\,\%$}{$\dfrac{1}{8}$}{$0{,}3$}
  \item Une baisse de $20\,\%$ est compensée par une hausse de :
        \qcm{$20\,\%$}{$25\,\%$}{$80\,\%$}{$120\,\%$}
  \item La droite qui passe par $A(2\,;1)$ et $B(4\,;7)$ a pour coefficient
        directeur :
        \qcm{$3$}{$\dfrac{1}{3}$}{$2$}{$-3$}
  \item La fonction $f(x)=3x-4$ est :
        \par a) croissante \quad b) décroissante \quad c) constante
  \item L'antécédent de $13$ par $f(x)=2x+3$ est :
        \qcm{$29$}{$5$}{$8$}{$6{,}5$}
  \item Le point d'abscisse $3$ de la droite d'équation $y=-2x+4$ a pour ordonnée :
        \qcm{$-2$}{$10$}{$2$}{$-10$}
  \item Une distance de $x$ miles vaut $1{,}6x$ kilomètres. Pour $x=25$, cela fait :
        \qcm{$40$~km}{$26{,}6$~km}{$15{,}6$~km}{$4$~km}
\end{enumerate}

\exo[Les inscrits d'un club de judo]
Une feuille de calcul donne le nombre d'inscrits d'un club de judo.
\begin{center}
\renewcommand{\arraystretch}{1.25}
\small
\begin{tabular}{|c|l|c|c|c|c|}
\hline
 & \multicolumn{1}{c|}{A} & B & C & D & E\\
\hline
1 & Année & 2020 & 2021 & 2022 & 2023\\
\hline
2 & Rang $x$ & $0$ & $1$ & $2$ & $3$\\
\hline
3 & Inscrits & $180$ & $171$ & $166$ & $156$\\
\hline
4 & Taux d'évolution & \large$\times$ & & & \\
\hline
\end{tabular}
\end{center}
\begin{enumerate}[label=\arabic*.]
  \item Calculer le taux d'évolution du nombre d'inscrits entre 2020 et 2021.
  \item Quelle formule saisir en C4 pour obtenir, par recopie vers la droite, les
        taux d'évolution d'une année sur l'autre ?
  \item On modélise le nombre d'inscrits par $f(x)=-8x+180$, où $x$ est le rang de
        l'année. Donner le sens de variation de $f$ et interpréter $-8$.
  \item Selon ce modèle, combien d'inscrits y aura-t-il en 2028 ?
  \item À partir de quelle année y aura-t-il moins de $100$ inscrits ?
\end{enumerate}

\exo[Les soldes]
Un magasin fait une remise de $25\,\%$ sur un casque à $80$~€ et de $10\,\%$ sur une
enceinte à $120$~€.
\begin{enumerate}[label=\arabic*.]
  \item Calculer le prix soldé de chaque article.
  \item On achète les deux articles. Calculer le prix total payé, puis le pourcentage
        global de remise.
  \item Après les soldes, le prix de l'enceinte augmente de $10\,\%$. Revient-il à
        $120$~€ ?
\end{enumerate}

\end{multicols}
\end{document}
